% -----------------------------------------------------------------------------
\section{Final System Architecture and Design}
\label{sec:architecture}
% -----------------------------------------------------------------------------

The final MoonAI architecture follows a GPU-first design. The CPU loads experiments and
settings, starts the UI, orchestrates kernel launches, and writes compact exports. The GPU owns
the simulation state and performs the dominant computation for sensing, inference, movement,
combat, reproduction, and metrics reduction.

\subsection{Architecture Overview}

\begin{figure}[H]
  \centering
  \includegraphics[width=\textwidth,height=0.62\textheight,keepaspectratio]{system_overview}
  \caption{High-level architecture of the final MoonAI runtime and analysis workflow.}
  \label{fig:system-overview}
\end{figure}

As shown in Figure~\ref{fig:system-overview}, the runtime begins with two external inputs:
\texttt{experiments.lua} and \texttt{settings.json}. The executable loads those files,
initializes the desktop application, creates or replaces run sessions, and coordinates GPU
simulation plus output logging. Completed run artifacts are then consumed by the Python
analysis pipeline to produce a self-contained HTML report.

\subsection{Module Decomposition}

The final repository is organized around clear module responsibilities.

\begin{table}[H]
  \centering
  \caption{Main implementation modules in the final codebase}
  \begin{tabularx}{\textwidth}{@{}p{3.4cm}X@{}}
    \toprule
    Module & Responsibility \\
    \midrule
    \texttt{main.rs} & Bootstraps the executable, resolves the runtime directory, loads
    settings and experiments, and starts the UI application. \\
    \texttt{experiment.rs} & Defines \texttt{SimulationConfig}, loads
    \texttt{experiments.lua}, injects defaults into Lua, and validates experiment values. \\
    \texttt{settings.rs} & Loads and saves \texttt{settings.json} and defines the persistent UI
    configuration surface. \\
    \texttt{src/ui/} & Implements the desktop application shell, run queue, active run session,
    render helpers, UI state, and custom world renderer. \\
    \texttt{src/sim/} & Defines the host-side simulation API, GPU readback structures, and CUDA
    bindings for the simulation and evolution kernels. \\
    \texttt{metrics.rs} & Writes structured runtime artifacts to the experiment output
    directory. \\
    \texttt{analysis/} & Discovers completed runs, groups conditions, generates plots, and
    renders a self-contained HTML analysis report. \\
    \bottomrule
  \end{tabularx}
\end{table}

This decomposition keeps configuration, runtime execution, rendering, metrics export, and
analysis separate enough to remain understandable while still functioning as one coherent
system.

\subsection{GPU-First Tick Flow}

\begin{figure}[H]
  \centering
  \includegraphics[width=0.92\textwidth,height=0.70\textheight,keepaspectratio]{tick_flow}
  \caption{Simplified tick flow for the final GPU-first simulation runtime.}
  \label{fig:tick-flow}
\end{figure}

The internal tick execution path is summarized in Figure~\ref{fig:tick-flow}. At each step,
the runtime builds spatial grids, computes sensor inputs, executes neural inference, updates
agent state, resolves interactions, performs reproduction, advances counters, and optionally
refreshes metrics and render snapshots.

Two cadences are intentionally separated in the design:

\begin{itemize}[nosep]
  \item \textbf{\texttt{report\_interval\_ticks}} controls when export artifacts such as
    \texttt{stats.csv}, \texttt{species.csv}, and \texttt{genomes.json} are refreshed.
  \item \textbf{UI speed multiplier} controls how often the interactive interface refreshes the
    visible world state.
\end{itemize}

This separation matters because reporting and rendering serve different purposes. Export should
remain consistent and comparable across runs, while UI refresh rate should remain adjustable
for live observation.

\subsection{Design Rationale}

The main architectural decisions were driven by three principles:

\begin{enumerate}[nosep]
  \item \textbf{Keep the hot path on the GPU.} Large populations make CPU-side simulation
    unattractive. GPU ownership of simulation state reduces host-side duplication and keeps the
    main path compact \cite{cuda}.
  \item \textbf{Expose the system through a usable interface.} The project is not only a kernel
    benchmark. It is also a tool for selecting experiments, watching live behavior, inspecting
    agents, and exporting results through a desktop application built with \texttt{egui} and
    \texttt{wgpu} \cite{egui,wgpu}.
  \item \textbf{Make every run analyzable after execution.} Structured metrics, species summaries,
    and representative genomes are part of the product, not optional extras.
\end{enumerate}

This combination of GPU execution, interactive runtime control, and post-run analysis is the
core of the final MoonAI design.
